% ============================================================
%  圆环设计说明（ID160 / OD200 / 环宽20 / 厚2）
%  编译：在本目录运行  xelatex 圆环设计说明.tex
% ============================================================
\documentclass[a4paper, 11pt]{ctexart}
\usepackage{geometry}
\geometry{left=2cm, right=2cm, top=2cm, bottom=2cm}
\usepackage{tikz}
\usepackage{booktabs}
\usepackage{amssymb}
\usetikzlibrary{arrows.meta}
\begin{document}

\begin{center}
{\LARGE\bfseries 圆环零件设计说明}\\[4pt]
{\normalsize 内径 160\,mm ｜ 环宽 20\,mm ｜ 外径 200\,mm ｜ 厚度 2\,mm}
\end{center}

\section*{一、尺寸定义}
\begin{center}
\renewcommand{\arraystretch}{1.3}
\begin{tabular}{lll}
\toprule
参数 & 数值 & 说明 \\
\midrule
内径 ID & 160\,mm & 给定 \\
环宽（径向） & 20\,mm & 给定 \\
\textbf{外径 OD} & \textbf{200\,mm} & 计算：160 + 2$\times$20 = 200 \\
厚度（轴向） & 2\,mm & 给定 \\
\bottomrule
\end{tabular}
\end{center}

\noindent 说明：按"环宽为径向宽度"理解，外径 $=$ 内径 $+$ 2\,$\times$\,环宽 $=$ 160\,$+$40 $=$ 200\,mm。

\section*{二、结构简图}

\begin{center}
\begin{tikzpicture}[x=0.5mm, y=0.5mm, >=Stealth]
  % 俯视图
  \begin{scope}[shift={(0,0)}]
    \draw[thick] (0,0) circle (100);
    \draw[thick] (0,0) circle (80);
    \draw[dashed] (-110,0) -- (110,0);
    \draw[dashed] (0,-110) -- (0,110);
    \draw[<->] (-100,0) -- (100,0);
    \draw[<->] (0,-80) -- (0,80);
    \node[below] at (0,-70) {$\varnothing$160};
    \node[above] at (0,20) {$\varnothing$200};
    \node at (0,-125) {俯视图};
  \end{scope}
  % 剖视图
  \begin{scope}[shift={(150,0)}, x=1.6mm, y=1.6mm]
    \draw[thick] (0,0) rectangle (20,2);
    \draw[<->] (0,2.8) -- (20,2.8);
    \node at (10,2.8) [above] {20（环宽）};
    \draw[<->] (22,0) -- (22,2);
    \node[below] at (22,1) {2};
    \node at (10,-1.5) {剖视图};
  \end{scope}
\end{tikzpicture}
\end{center}

\section*{三、质量估算}
体积 $V=\dfrac{\pi}{4}\left(200^{2}-160^{2}\right)\times 2 \approx 22\,619\ \mathrm{mm^{3}} \approx 22.6\ \mathrm{cm^{3}}$（SolidWorks 实测 22\,619.5\,mm³）。

\begin{center}
\begin{tabular}{lcc}
\toprule
材料 & 密度 (g/cm³) & 估算质量 (g) \\
\midrule
钢 & 7.85 & $\approx$178 \\
铝 & 2.70 & $\approx$61 \\
铜 & 8.96 & $\approx$203 \\
钛合金 & 4.51 & $\approx$102 \\
氧化铝陶瓷 & 3.90 & $\approx$88 \\
\bottomrule
\end{tabular}
\end{center}

\section*{四、交付文件}
\begin{center}
\begin{tabular}{ll}
\toprule
文件 & 用途 \\
\midrule
\texttt{Ring\_ID160\_OD200\_T2.SLDPRT} & SolidWorks 源文件（可改尺寸） \\
\texttt{Ring\_ID160\_OD200\_T2.step} & 通用三维格式（导入其他 CAD） \\
\texttt{Ring\_ID160\_OD200\_T2.stl} & 3D 打印 \\
\texttt{Ring\_iso.png / Ring\_top.png} & 预览图 \\
\bottomrule
\end{tabular}
\end{center}

\noindent\textbf{备注}：本次未标注公差；若用于配合件（如轴孔配合），需要给出配合公差与形位公差（同轴度、平面度），或说明使用场景（垫圈/夹具/炉具垫环等），我再补标注。

\section*{五、理解确认}
若"宽 20\,mm"指的是\textbf{轴向高度}、"厚度 2\,mm"指\textbf{径向壁厚}，则外径应为 $160+2\times2=164$\,mm，结构变为高 20\,mm 的薄壁圆筒。本图按前者（环宽为径向）制作；如需按后者，可立即重做。

\end{document}
